% Options for packages loaded elsewhere
% Options for packages loaded elsewhere
\PassOptionsToPackage{unicode}{hyperref}
\PassOptionsToPackage{hyphens}{url}
\PassOptionsToPackage{dvipsnames,svgnames,x11names}{xcolor}
%
\documentclass[
  letterpaper,
  DIV=11,
  numbers=noendperiod]{scrartcl}
\usepackage{xcolor}
\usepackage[margin=1in]{geometry}
\usepackage{amsmath,amssymb}
\setcounter{secnumdepth}{-\maxdimen} % remove section numbering
\usepackage{iftex}
\ifPDFTeX
  \usepackage[T1]{fontenc}
  \usepackage[utf8]{inputenc}
  \usepackage{textcomp} % provide euro and other symbols
\else % if luatex or xetex
  \usepackage{unicode-math} % this also loads fontspec
  \defaultfontfeatures{Scale=MatchLowercase}
  \defaultfontfeatures[\rmfamily]{Ligatures=TeX,Scale=1}
\fi
\usepackage{lmodern}
\ifPDFTeX\else
  % xetex/luatex font selection
\fi
% Use upquote if available, for straight quotes in verbatim environments
\IfFileExists{upquote.sty}{\usepackage{upquote}}{}
\IfFileExists{microtype.sty}{% use microtype if available
  \usepackage[]{microtype}
  \UseMicrotypeSet[protrusion]{basicmath} % disable protrusion for tt fonts
}{}
\makeatletter
\@ifundefined{KOMAClassName}{% if non-KOMA class
  \IfFileExists{parskip.sty}{%
    \usepackage{parskip}
  }{% else
    \setlength{\parindent}{0pt}
    \setlength{\parskip}{6pt plus 2pt minus 1pt}}
}{% if KOMA class
  \KOMAoptions{parskip=half}}
\makeatother
% Make \paragraph and \subparagraph free-standing
\makeatletter
\ifx\paragraph\undefined\else
  \let\oldparagraph\paragraph
  \renewcommand{\paragraph}{
    \@ifstar
      \xxxParagraphStar
      \xxxParagraphNoStar
  }
  \newcommand{\xxxParagraphStar}[1]{\oldparagraph*{#1}\mbox{}}
  \newcommand{\xxxParagraphNoStar}[1]{\oldparagraph{#1}\mbox{}}
\fi
\ifx\subparagraph\undefined\else
  \let\oldsubparagraph\subparagraph
  \renewcommand{\subparagraph}{
    \@ifstar
      \xxxSubParagraphStar
      \xxxSubParagraphNoStar
  }
  \newcommand{\xxxSubParagraphStar}[1]{\oldsubparagraph*{#1}\mbox{}}
  \newcommand{\xxxSubParagraphNoStar}[1]{\oldsubparagraph{#1}\mbox{}}
\fi
\makeatother


\usepackage{longtable,booktabs,array}
\usepackage{calc} % for calculating minipage widths
% Correct order of tables after \paragraph or \subparagraph
\usepackage{etoolbox}
\makeatletter
\patchcmd\longtable{\par}{\if@noskipsec\mbox{}\fi\par}{}{}
\makeatother
% Allow footnotes in longtable head/foot
\IfFileExists{footnotehyper.sty}{\usepackage{footnotehyper}}{\usepackage{footnote}}
\makesavenoteenv{longtable}
\usepackage{graphicx}
\makeatletter
\newsavebox\pandoc@box
\newcommand*\pandocbounded[1]{% scales image to fit in text height/width
  \sbox\pandoc@box{#1}%
  \Gscale@div\@tempa{\textheight}{\dimexpr\ht\pandoc@box+\dp\pandoc@box\relax}%
  \Gscale@div\@tempb{\linewidth}{\wd\pandoc@box}%
  \ifdim\@tempb\p@<\@tempa\p@\let\@tempa\@tempb\fi% select the smaller of both
  \ifdim\@tempa\p@<\p@\scalebox{\@tempa}{\usebox\pandoc@box}%
  \else\usebox{\pandoc@box}%
  \fi%
}
% Set default figure placement to htbp
\def\fps@figure{htbp}
\makeatother





\setlength{\emergencystretch}{3em} % prevent overfull lines

\providecommand{\tightlist}{%
  \setlength{\itemsep}{0pt}\setlength{\parskip}{0pt}}



 



\usepackage{tcolorbox}
% シンプルな枠線の定義
\newtcolorbox{reviewerbox}{
  % sharp corners,                % 角を丸めない
  colback=white,                % 背景は白
  colframe=black,               % 枠線は黒
  % boxrule=0.5pt,                % 線の太さ
  fontupper=\itshape,           % 中身をイタリックにする
  before skip=12pt,             % 前後の余白
  after skip=12pt
}

\usepackage{longtable}
\usepackage{booktabs}
\usepackage{array}
\KOMAoption{captions}{tableheading}
\makeatletter
\@ifpackageloaded{caption}{}{\usepackage{caption}}
\AtBeginDocument{%
\ifdefined\contentsname
  \renewcommand*\contentsname{Table of contents}
\else
  \newcommand\contentsname{Table of contents}
\fi
\ifdefined\listfigurename
  \renewcommand*\listfigurename{List of Figures}
\else
  \newcommand\listfigurename{List of Figures}
\fi
\ifdefined\listtablename
  \renewcommand*\listtablename{List of Tables}
\else
  \newcommand\listtablename{List of Tables}
\fi
\ifdefined\figurename
  \renewcommand*\figurename{Figure}
\else
  \newcommand\figurename{Figure}
\fi
\ifdefined\tablename
  \renewcommand*\tablename{Table}
\else
  \newcommand\tablename{Table}
\fi
}
\@ifpackageloaded{float}{}{\usepackage{float}}
\floatstyle{ruled}
\@ifundefined{c@chapter}{\newfloat{codelisting}{h}{lop}}{\newfloat{codelisting}{h}{lop}[chapter]}
\floatname{codelisting}{Listing}
\newcommand*\listoflistings{\listof{codelisting}{List of Listings}}
\makeatother
\makeatletter
\makeatother
\makeatletter
\@ifpackageloaded{caption}{}{\usepackage{caption}}
\@ifpackageloaded{subcaption}{}{\usepackage{subcaption}}
\makeatother
\usepackage{bookmark}
\IfFileExists{xurl.sty}{\usepackage{xurl}}{} % add URL line breaks if available
\urlstyle{same}
\hypersetup{
  pdftitle={Response to Reviewers},
  colorlinks=true,
  linkcolor={blue},
  filecolor={Maroon},
  citecolor={Blue},
  urlcolor={Blue},
  pdfcreator={LaTeX via pandoc}}


\title{Response to Reviewers}
\author{}
\date{}
\begin{document}
\maketitle


\section{Response to Reviewer 1}\label{response-to-reviewer-1}

This sample text was generated by Google Gemini.
We would like to thank Reviewer 1 for their rigorous and constructive feedback on our manuscript, ``The History of Animagus: From Medieval Transfiguration to Modern Registration.''
We have addressed each point as detailed below.

\subsection{Comment \#1}\label{comment-1}

\begin{reviewerbox}

The authors state that the Animagus Registry is an exhaustive record of all practitioners in Great Britain. However, historical evidence from the 1970s suggests the existence of several unregistered Animagi at Hogwarts. The manuscript should acknowledge these non-compliant cases to maintain historical accuracy.

\end{reviewerbox}

Thank you for this crucial observation. We agree that the distinction between registered and unregistered Animagi is vital for a comprehensive history. We have updated the introduction and the ``Ministry Oversight'' section to include the notable cases of the `Marauders' cohort.

\begin{longtable*}{p{0.18\textwidth} p{0.36\textwidth} p{0.36\textwidth}}
                    \toprule
                    \textbf{Location} & \textbf{Before} & \textbf{After} \\
                    \midrule
                    \endhead
                    Page 3, Line 45 & The Ministry of Magic maintains a complete and infallible registry of
all Animagi, ensuring every transformation is monitored for public
safety. & While the Ministry of Magic maintains an official registry, historical
records---including the `Marauders' incident---reveal that
\textbf{unregistered transformations have occurred}, often bypassing
Ministry oversight for years. \\
\midrule
\end{longtable*}

\subsection{Comment \#2}\label{comment-2}

\begin{reviewerbox}

The description of the transformation process is somewhat vague. Specifically, the requirement to keep a Mandrake leaf in one's mouth for an entire month is a defining (and grueling) feature of the process and deserves explicit mention.

\end{reviewerbox}

We appreciate this suggestion. We have added a new subsection, ``The Mandrake Protocol,'' providing a detailed account of the physiological and psychological challenges of the transformation process.

\begin{longtable*}{p{0.18\textwidth} p{0.36\textwidth} p{0.36\textwidth}}
                    \toprule
                    \textbf{Location} & \textbf{Before} & \textbf{After} \\
                    \midrule
                    \endhead
                    Page 7, Section 2.3 & The transformation involves complex potion-making and high-level
incantations over an extended period. & The transformation begins with the \textbf{Mandrake Protocol}, requiring
the practitioner to keep a single Mandrake leaf in their mouth from one
full moon to the next, followed by the consumption of an Animagus Potion
during a lightning storm. \\
\midrule
Page 12, Conclusion & In conclusion, the history of Animagus shows a trend toward stricter
regulation. & In conclusion, the history of Animagus reflects a constant tension
between \textbf{Ministry regulation and the inherent secrecy} of those
who achieve this rare magical feat. \\
\midrule
\end{longtable*}

\newpage{}

\section{Response to Reviewer 2}\label{response-to-reviewer-2}

Reviewer 2 raised interesting points regarding the psychological impact of long-term animal form maintenance.

\subsection{Comment \#1}\label{comment-1-1}

\begin{reviewerbox}

The case of Peter Pettigrew (who remained in rat form for twelve years) provides a unique data point on the `human-animal cognitive overlap.' The authors should discuss whether prolonged transformation leads to a permanent loss of human identity.

\end{reviewerbox}

This is a fascinating point that aligns with our secondary research interest in cognitive neuroscience. We have added a paragraph discussing the ``Pettigrew Effect,'' where long-term transformation may result in a significant shift toward animalistic behavioral patterns.

\begin{longtable*}{p{0.18\textwidth} p{0.36\textwidth} p{0.36\textwidth}}
                    \toprule
                    \textbf{Location} & \textbf{Before} & \textbf{After} \\
                    \midrule
                    \endhead
                    Page 15, Discussion & Pettigrew's case is primarily seen as a failure of Ministry detection. & Pettigrew's twelve-year residence as a common garden rat suggests a
potential \textbf{`identity erosion'}, where human cognitive structures
are increasingly dominated by the instincts of the animal form. \\
\midrule
\end{longtable*}




\end{document}
